\documentclass[11pt]{beamer}
\usepackage[utf8]{inputenc}
\usepackage{amsmath, amssymb, graphicx, listings, xcolor}

% Presentation Theme
\usetheme{Madrid}
\usecolortheme{default}

% Code Block Styling
\lstset{
    language=Python,
    basicstyle=\ttfamily\scriptsize,
    keywordstyle=\color{blue}\bfseries,
    commentstyle=\color{green!50!black},
    frame=single,
    backgroundcolor=\color{gray!10}
}

\title[Continuous-Thrust Trajectories]{Astrodynamics 2026-1}
\subtitle{Finite Maneuvers \& Low-Thrust Spirals}
\author{Prof. Juan}
\institute{Universidad de Antioquia - Aerospace Engineering}
\date{Spring 2026}

\begin{document}

% --- Title Slide ---
\begin{frame}
    \titlepage
\end{frame}

% ==========================================
% SECTION 1: THE PARADIGM SHIFT
% ==========================================
\section{The Limits of Impulsive Maneuvers}

\begin{frame}{The Impulsive Assumption}
    Until now, we have modeled engine burns as \textbf{Impulsive Maneuvers}. 
    
    \vspace{0.3cm}
    \textbf{The Assumption:} 
    \begin{itemize}
        \item The burn time ($\Delta t$) is infinitesimally small.
        \item The position vector ($\boldsymbol{r}$) does not change during the burn.
        \item The velocity vector changes instantly: $\boldsymbol{v}^+ = \boldsymbol{v}^- + \Delta\boldsymbol{v}$.
    \end{itemize}
    
    \vspace{0.3cm}
    \textbf{The Reality:}
    Real engines have finite thrust. A chemical engine might burn for 5 minutes. An ion thruster (Solar Electric Propulsion) might burn for \textbf{6 months}! We can no longer assume position is constant during the maneuver.
\end{frame}

% ==========================================
% SECTION 2: GMAT MATHEMATICAL SPECS
% ==========================================
\section{Equations of Motion}

\begin{frame}{Finite Maneuvers: Equations of Motion}
    According to the GMAT Mathematical Specifications, when a \texttt{FiniteBurn} is active, the spacecraft's state vector expands. We must continuously integrate acceleration and mass flow.
    
    \vspace{0.3cm}
    \textbf{1. The Acceleration Vector:}
    \begin{equation}
        \ddot{\boldsymbol{r}} = -\frac{\mu}{r^3}\boldsymbol{r} + \boldsymbol{a}_{perturbations} + \frac{\boldsymbol{T}}{m(t)}
    \end{equation}
    Where $\boldsymbol{T}$ is the thrust vector and $m(t)$ is the continuously decreasing mass.
    
    \vspace{0.3cm}
    \textbf{2. The Mass Depletion Rate:}
    \begin{equation}
        \dot{m} = -\frac{|\boldsymbol{T}|}{I_{sp} g_0}
    \end{equation}
    Because mass is in the denominator of acceleration, a constant thrust force $\boldsymbol{T}$ results in an \textit{increasing} acceleration over time!
\end{frame}

% ==========================================
% SECTION 3: VALLADO'S SPIRAL
% ==========================================
\section{Low-Thrust Spirals}

\begin{frame}{Vallado: The Quasi-Circular Spiral}
       
    
    \begin{equation}
        \Delta V_{spiral} \approx |v_{c1} - v_{c2}| = \left| \sqrt{\frac{\mu}{r_1}} - \sqrt{\frac{\mu}{r_2}} \right|
    \end{equation}
    
    \vspace{0.2cm}
    \textit{Note: This is simply the absolute difference between the circular velocities of the two orbits.}
\end{frame}

\begin{frame}{The Great Trade-Off: Mass vs. Time}
    \textbf{The Gravity Penalty:}
    \begin{equation}
        \Delta V_{spiral} > \Delta V_{Hohmann}
    \end{equation}
    Low thrust suffers from "gravity losses" because we spend much of the transfer thrusting inefficiently (not at periapsis).
    
    \vspace{0.3cm}
    \textbf{Why do we use it?}
    Electric thrusters have Specific Impulses ($I_{sp}$) of 3,000 to 5,000 seconds (10x better than chemical rockets). Using the Rocket Equation:
    \begin{equation}
        m_f = m_0 e^{-\frac{\Delta V}{I_{sp} g_0}}
    \end{equation}
    Even though $\Delta V$ is higher, the massive $I_{sp}$ means we deliver significantly more payload mass. We trade \textbf{Time} (months instead of days) for \textbf{Mass}.
\end{frame}

% ==========================================
% SECTION 4: GMAT IMPLEMENTATION
% ==========================================
\section{GMAT Software Implementation}

\begin{frame}[fragile]{Modeling Finite Burns in GMAT}
    To model this in GMAT, we must abandon the \texttt{ImpulsiveBurn} object and construct the physical hardware.
    
\begin{lstlisting}
% 1. Create the Hardware
Create ElectricThruster IonEngine;
IonEngine.DecrementMass = true;
IonEngine.Isp = 3000;

Create ElectricTank XenonTank;
Jupiter_SC.Thrusters = {IonEngine};
Jupiter_SC.Tanks = {XenonTank};

% 2. Execute the Finite Burn (Inside Mission Sequence)
BeginFiniteBurn IonEngine(Jupiter_SC);
    
    % Propagate until target mass or SMA is reached
    Propagate DefaultProp(Jupiter_SC) {Jupiter_SC.SMA = 42164};
    
EndFiniteBurn IonEngine(Jupiter_SC);
\end{lstlisting}
\end{frame}

\end{document}